\section{Yêu cầu của bài tập lớn}
\label{sec:1_4}
Theo đề bài môn Học máy, sản phẩm cần thể hiện đầy đủ quy trình xử lý và thực nghiệm trên loại dữ liệu đã chọn~\cite{ml-assignment}. Với bài dữ liệu bảng UCI Adult, các yêu cầu được tổ chức như sau:

\begin{itemize}
  \item \textbf{Khảo sát dữ liệu:} thực hiện EDA, mô tả các thuộc tính, phân phối nhãn, giá trị thiếu và những vấn đề ảnh hưởng đến việc xây dựng mô hình.
  \item \textbf{Pipeline có thể cấu hình:} cho phép lựa chọn và điều chỉnh kỹ thuật tiền xử lý, giảm chiều và mô hình; sử dụng dữ liệu có thuộc tính phân loại và giá trị thiếu để thực hành encoding và imputation.
  \item \textbf{Thí nghiệm và đánh giá:} so sánh các cấu hình, báo cáo các chỉ số phù hợp và giải thích ưu điểm, hạn chế của kết quả thu được.
  \item \textbf{Notebook thực thi:} toàn bộ notebook phải chạy được bằng chức năng \texttt{Run all}; các bước chuẩn bị và tải dữ liệu từ nguồn công khai được thực hiện tự động trong môi trường Colab.
  \item \textbf{Tổ chức sản phẩm:} cung cấp báo cáo PDF, notebook, các module hỗ trợ khi có và các tệp đặc trưng định dạng \texttt{.npy} hoặc \texttt{.h5}; tổ chức rõ các thư mục \texttt{notebooks/}, \texttt{modules/}, \texttt{reports/} và \texttt{features/}.
  \item \textbf{Làm việc nhóm:} trình bày phân công, tỷ lệ đóng góp và minh chứng hoạt động của các thành viên trong báo cáo; thực hiện báo cáo tiến độ theo yêu cầu môn học.
\end{itemize}

Pipeline truyền thống là phần cốt lõi của bài làm. Việc xây dựng thêm MLP và tổ chức mã nguồn trên GitHub được đưa vào kế hoạch mở rộng, gắn với yêu cầu so sánh phương pháp và khả năng tái lập.
